%$Id: config_nux.tex 145 2005-03-25 08:26:35Z myk $

\clearpage
\pagebreak

\bghdr{images/fond-nux}

%\begin{center}
%\includegraphics{images/logo_Linux}
%\end{center}

\subsection{Configuration sous Linux}

\subsubsection{Configuration IP}
Quelle que soit ta distribution, il faut que tu d\'ecides d'un nom de machine --- ton pseudo en g\'en\'eral.
Tu as \'egalement besoin de conna�tre les informations suivantes :
\begin{description}
  \item [ton adresse IP] not\'ee dor\'enavant \server{129.104.AAA.BBB}
  \item [l'IP de ta passerelle (\emph{gateway})] not\'ee \server{129.104.GGG.CCC}
  \item [ton masque de sous-r\'eseau] \server{255.255.FFF.DDD}
  \item [ton broadcast] qui est de la forme \server{129.104.GGG.EEE} avec $EEE = 255 - DDD$.
        \textcolor{red}{(attention, $GGG = AAA$ n'est valable que dans les nouveaux caserts, le PEM et le BEM. Pour Foch, Fayolle et Maunoury ce n'est donc pas n\'ecessairement vrai que $GGG = AAA$.)}
\end{description}

Toutes les informations n\'ecessaires se trouvent page \pageref{calcul_ip}.

Bien s\^ur, l'ensemble des manipulations doit se faire en tant que root.

\paragraph{\distrib{Sous Mandriva}}
{\bf Configuration du r\'eseau}
Il existe dans cette distribution une interface graphique pour configurer le r\'eseau, bien qu'il soit possible,
comme sur tout syst\`eme linux, de le faire \`a la main \`a partir des fichiers de configuration dans \rep{/etc/}.

Lance \app{drakconf} en tant que root, puis clique sur l'onglet \menu{R\'eseau \& Internet} et enfin sur \menu{Configurer une nouvelle connexion Internet (LAN, ISDN, ADSL, ...)}.
Choisis ensuite \menu{Connexion � travers un r\'eseau local(LAN)}, puis la carte r\'eseau sur laquelle tu as branch\'e ton c\^able r\'eseau.
Tu as autant de \texttt{eth0}, \texttt{eth1}, \ldots que de cartes r\'eseau d\'etect\'ees.

Sur la fen\^etre suivante, choisis l'option \menu{configuration manuelle}.  La configuration du r\'eseau commence alors. 
Dans la premi�re fen\^etre, entre :
\begin{description}
       \item[ton adresse IP] entre celle que tu as calcul\'ee : \server{129.104.AAA.BBB} (cf. page \pageref{calcul_ip})
       \item[ton masque de sous-r\'eseau] : \server{255.255.FFF.DDD} (cf. page \pageref{calcul_ip})
\end{description}

Ensuite, tu dois encore entrer :

\begin{tabular}{@{\bf }l@{\xleaders\copy\point\ecart}l}
       Nom d'h�te (non obligatoire) & ton\_pseudo \\
       Serveur DNS 1       & \server{129.104.201.53} \\
       Serveur DNS 2       & \server{129.104.201.54} \\
       Serveur DNS 3       & \server{129.104.201.51} \\
       Domaine recherch\'e & eleves.polytechnique.fr polytechnique.fr \\
       Passerelle          & \server{129.104.GGG.CCC} (cf. page \pageref{calcul_ip}) \\
       P\'eriph\'erique passerelle  & aucun  \\
\end{tabular}

Clique ensuite deux fois sur suivant. (Tu n'as pas besoin d'entrer de nom d'h\^ote ZeroConf.)
L�, tu choisis l'option \menu{Oui}, puis tu valides et tu as normalement un message
qui t'annonce que ta configuration r\'eseau est termin\'ee.
Tu peux la tester en pinguant frankiz, qui doit donner quelque chose comme :

\cmdline{\$ ping frankiz\\
PING frankiz.eleves.polytechnique.fr (129.104.201.51) 56(84) bytes of data.\\
64 bytes from frankiz (129.104.201.51): ...
}


N'oublie pas d'ouvrir les ports qui vont bien (21 pour un serveur ftp,
5050, 5053 et 5055 pour qRezix, ...) dans ton firewall pour pouvoir pleinement profiter du r\'eseau.

{\bf Configuration du gestionnaire de paquets :}
Le BR poss\`ede un miroir PLF, que tu peux utiliser en choisissant
l'onglet \menu{Gestionnaire de logiciels} dans l'utilitaire \app{drakconf}
puis en cliquant sur \menu{ajouter la source personalis\'ee}, \`a choisir \menu{Serveur FTP}
et en mettant comme url \url{ftp://miroir/plf/mandrake/free/2006.0/}
pour le premier miroir et \url{ftp://miroir/plf/mandrake/non-free/2006.0/}
pour le second. Il te faut aussi sp\'ecifier le chemin vers le
hdlist.cz, en mettant dans les deux cas ./hdlist.cz dans le champ
\menu{Chemin relatif vers hdlist/synthesis}. Enfin, n'oublie pas de
d�sactiver le proxy pour les utiliser.


\paragraph{\distrib{Sous Gentoo}}

Si tu n'as jamais fait de configuration r\'eseau sur ta Gentoo,
il faudra certainement cr\'eer les fichiers qui suivent.

Le fichier \file{/etc/hostname} contient ton nom de machine.
Tu peux \'editer \file{/etc/hostname} avec ton \'editeur pr\'ef\'er\'e
(\app{vi} ou \app{vim}, \app{pico}, voire \app{emacs} si tu aimes),
mais la commande ci-dessous est suffisante :

\cmdline{echo 'ton\_pseudo' > /etc/hostname}

Le fichier \file{/etc/resolv.conf} d\'ecrit comment r\'esoudre les noms DNS.
La premi\`ere ligne donne le domaine sur lequel ta machine est, 
ensuite viennent les suffixes \`a utiliser par d\'efaut,
et les lignes suivantes indiquent les serveurs de noms,
ceux qui associent le nom \server{frankiz} ou \server{ton\_pseudo}
aux IP \server{129.104.201.51} ou \server{'ton\_ip'}.
Le fichier contient donc :

\noindent \cmdline{
domain eleves.polytechnique.fr\\
search eleves.polytechnique.fr polytechnique.fr\\
nameserver 129.104.201.53\\
namaserver 129.104.201.52}

Enfin, le fichier \file{/etc/conf.d/net} contient la description de l'interface r\'eseau et les informations de routage.
Il indique ton IP, ton sous-r\'eseau, ton masque de sous-r\'eseau et la passerelle pour sortir de ton sous-r\'eseau .
Bien s�r, remplaces \texttt{eth0} par le nom de l'interface r\'eseau que tu utilises,
si ce n'est pas la carte \textit{ethernet} par d\'efaut.

\noindent \cmdline{iface\_eth0="129.104.AAA.BBB broadcast 129.104.GGG.EEE $\backslash$\\
                            netmask 255.255.FFF.DDD"\\
                            gateway="eth0/129.104.GGG.CCC"}

Tu fais le grand test en rechargeant ta configuration r\'eseau par :

\cmdline{/etc/init.d/net.eth0 restart}

puis en pingant \fkz\ par exemple. Tu dois obtenir quelque chose comme :

\cmdline{root:~\# ping frankiz\\
PING frankiz (129.104.201.51) 56(84) bytes of data.\\
64 bytes from frankiz (129.104.201.51): ...}

Pour pouvoir utiliser \app{emerge} \`a travers le proxy de l'\'ecole,
il faut d\'efinir les variables d'environnement ci-dessous dans le fichier \file{/etc/make.conf} :

\cmdline{http\_proxy=http://kuzh.polytechnique.fr:8080\\
GENTOO\_MIRRORS="ftp://miroir/gentoo http://gentoo.osuosl.org/"\\
SYNC="rsync://rsync/gentoo-portage"}

Tu peux \'evidemment ajouter d'autres miroirs (s\'epar\'es par des espaces) dans ta liste
mais \url{ftp://miroir} \'etant interne, il sera toujours beaucoup plus rapide que les autres.
On te conseille ici \url{http://gentoo.osuosl.org/} qui est un miroir tr\`es complet mais tr\`es lent.

\input ubuntu.tex

\subsubsection{Configuration antivirus \footnotesize{(elle est dr\^ole celle-l\`a hein ?)}
}

\subsubsection{Configuration firewall}

La solution la plus simple pour se faire un firewall sous linux est d'utiliser les iptables.
Pour ceci la premi\`ere \'etape est d'installer le paquet \app{iptables} pour ta distribution.
Pour savoir comment configurer ton firewall pour le r\'eseau de l'X, consulte la FAQ de Frankiz.

\subsubsection{Configuration navigateur web}

\imagepos{images/nux_proxy_firefox}{0.60}{Configuration du proxy sous Firefox}{ht}

\flimage{images/nux_firefox_icon}{0.12}{l}
Le BR te conseille d'utiliser \app{Firefox}, ou \app{Konqueror} (le navigateur fourni par d\'efaut avec KDE).
Dans tous les cas, la seule configuration \`a mettre est celle du proxy.
Pour \app{Firefox}, il suffit d'aller dans \menu{Edit}, \menu{Preferences} et dans l'onglet \menu{General}
cliquer sur \menu{Connection Settings} ; ensuite tu coches la case
\menu{D\'etection automatique du proxy pour ce r\'eseau}, et c'est bon.

Sous \app{Konqueror}, cela se trouve dans le menu \menu{Configuration}, \menu{Configurer Konqueror},
dans l'onglet \menu{Serveur mandataire}.

% \vfill
\pagebreak

\subsubsection{Configuration mail}

\flimage{images/nux_kmail_icon}{0.12}{l}
Le client mail le plus utilis\'e est \app{Kmail}, mais il en existe bien s\^ur d'autres comme \app{Thunderbird}.
La configuration est semblable, quel que soit le client que tu utilises.

Va dans \menu{Configuration}, \menu{Configurer Kmail}.
Choisis la rubrique \menu{R\'eseau}.
Commence par cr\'eer un nouveau compte dans l'onglet \menu{R\'eception des messages}
en cliquant sur le bouton \menu{Ajouter\ldots} et choisis le type POP3.

\imagepos{images/nux_config_kmail_pop}{0.66}{Configuration de la r\'eception des messages sous Kmail}{pht}

Utilise les param\`etres suivants pour configurer l'onglet \menu{G\'en\'eral} :
\begin{description}
  \item[Nom] le nom du compte, par exemple : Mails Poly
  \item[Utilisateur] rentre le login \server{poly} que t'a fourni la DSI \`a ton arriv\'ee sur le plateau
  \item[Mot de passe] et l\`a le mot de passe \server{poly}
  \item[Serveur] \server{poly.polytechnique.fr}
  \item[Port] 995
\end{description}

Ensuite, va dans l'onglet \menu{Extras} et coche la case \menu{Utiliser SSL pour s\'ecuriser les t\'el\'echargements}.

\imagepos{images/nux_config_kmail_smtp}{0.55}{Configuration de l'envoi des messages sous Kmail}{pht}

Maintenant, dans l'onglet \menu{Envoi des messages} clique sur le bouton \menu{Ajouter\ldots}.
Utilise les param\`etres suivants pour le configurer :
\begin{description}
  \item[Nom] le m\^eme nom de compte que pr\'ec\'edemment
  \item[Serveur] \server{poly.polytechnique.fr}
  \item[Port] 25
\end{description}
Sinon, laisse toutes les cases d\'ecoch\'ees.

\subsubsection{Configuration news}

\flimage{images/nux_knode_icon}{0.12}{l}
Le client news le plus utilis\'e est \app{Knode}.
Parmi les autres clients news, citons \app{Thunderbird} ou \app{slrn}.
Ici aussi, la configuration est presque ind\'ependante du logiciel choisi.

Sous \app{Knode}, c'est dans le menu \menu{Configuration}, puis \menu{Configurer Knode}.
Va dans la rubrique \menu{Comptes, Forums de discussion} et cr\'ee un compte en cliquant sur \menu{Nouveau\ldots}.

\imagepos{images/nux_config_knode}{0.7}{Configuration de Knode}{ht}

Remplis l'onglet \menu{Serveur} avec les informations suivantes :
\begin{description}
  \item[Nom] ce que tu veux pour d\'ecrire ce compte, par exemple 'News Frankiz'
  \item[Serveur] \server{frankiz.polytechnique.fr}
  \item[Port] 119
\end{description}

Ensuite occupe-toi de l'onglet \menu{Identit\'e} :
\begin{description}
  \item[Nom] mets ton pseudo dans ce champ
  \item[Organisation] X, �cole Polytechnique, comme tu le sens
  \item[Adresse \'electronique] ton adresse mail, pour que les gens puissent te r\'epondre par mail.
\end{description}

Enfin, pour que \app{Knode} puisse envoyer des mails, il faut aller dans la rubrique \menu{Comptes},
sous-rubrique \menu{Courrier \'electronique}, et choisir comme serveur d'envoi de mails \server{poly.polytechnique.fr},
port 25 --- c'est exactement la m\^eme configuration SMTP que \app{Kmail}.

Si tu veux mettre une signature \`a la fin des messages que tu posteras,
il te suffit de la mettre dans l'onglet \menu{Identit\'e}.
Sur la plupart des clients la signature est interpr\'et\'ee comme ext\'erieure au message
et n'est en particulier pas incluse dans le texte cit\'e lorsque tu r\'eponds \`a un message.
Pour d\'efinir une signature \`a la main, il suffit de mettre \verb*+-- +\ (c'est \`a dire -{}-<espace>)
sur une ligne, et tout ce qui suivra cette ligne composera ta signature.

Il ne te reste plus qu'\`a t'inscrire \`a des newsgroups
(reporte-toi \`a la page \pageref{newsgroups} pour plus d'infos) et \`a poster !

Pour te connecter aux serveurs de news de Polytechnique.org, qui ont un acc\`es s\'ecuris\'e,
avec \app{Knode}, il y a une petite subtilit\'e car il ne g\`ere pas le SSL.
Il faut installer \app{stunnel} qui permet de d\'efinir une redirection SSL de port.
Dans \file{/etc/stunnel.conf} ou \file{/etc/stunnel/stunnel.conf} selon ta distribution,
mets les lignes suivantes (les trois premi\`eres y sont en principe d\'ej\`a) :

\cmdline{\# location of pid file\\
pid = /etc/stunnel/stunnel.pid\\
\\
\# user to run as\\
setuid = stunnel\\
setgid = stunnel\\
\\
\# Use it for client mode\\
client = yes\\
\\
\# sample service-level configuration\\
\\
{[}nntps{]}\\
accept  = 1119\\
connect = ssl.polytechnique.org:563\\
TIMEOUTclose = 0
}

Il ne te reste plus qu'\`a lancer \app{stunnel} par :

\cmdline{/etc/init.d/stunnel start}

Et tu peux ainsi lire les news de Polytechnique.org en mettant \server{localhost} comme serveur
et \server{1119} comme port.
Il faut aussi que tu coches \menu{Le serveur exige une identification}
et que tu rentres ton nom d'utilisateur \`a Polytechnique.org et ton mot de passe,
que tu peux d\'efinir sur \url{http://www.polytechnique.org/acces\_smtp.php}.
